\chapter{Appearance and Representation}
\label{ch:appearance}

\section{The Representation panel}

\subsection{Representation mode}

\ui{Mode} selects how atoms and bonds are drawn:

\begin{description}
  \item[\ui{Ball-and-Stick}] Spheres at a fraction of the covalent radius,
    bonds as cylinders. The default.
  \item[\ui{Space-filling (CPK)}] Spheres at approximately the van der
    Waals radius; bonds are hidden inside the atoms.
  \item[\ui{Wireframe}] Bonds only, drawn as lines with points at the atom
    positions --- the cheapest mode for very large systems.
\end{description}

\ui{Atom radius} and \ui{Bond width} scale everything globally, 0.20--3.00,
each with a slider and a synchronised spin box for exact values.

\subsection{Colouring atoms}
\label{sec:color-modes}

\ui{Color by} offers four mappings, applied consistently to atoms and to
their halves of each bond:

\begin{description}
  \item[\ui{Element (CPK)}] The Jmol palette, with per-element overrides.
  \item[\ui{Coordination number (CN)}] Discrete coordination numbers on a
    gradient.
  \item[\ui{Generalized CN (GCN)}] Continuous generalised coordination
    numbers --- excellent for distinguishing terraces, steps, edges and
    vertices on nanoparticles and slabs.
  \item[\ui{Custom property}] Any per-atom scalar field the structure
    carries: charges, force magnitudes, extended-XYZ columns, or a computed
    field such as \texttt{local\_entropy}. \ui{Property} selects which.
\end{description}

For the three scalar modes, \ui{Gradient} chooses the colour map ---
\ui{Viridis}, \ui{Plasma}, \ui{Turbo}, \ui{Inferno}, \ui{Magma},
\ui{Cividis}, \ui{Hot}, \ui{Afmhot}, \ui{Coolwarm}, \ui{Rainbow} --- and
\ui{Invert palette} reverses the scalar-to-colour mapping, matching
matplotlib's \texttt{\_r} variants. \ui{Range} reads back the min and max
of the active field, so the legend is never ambiguous.

\begin{caltip}
Viridis and Cividis are perceptually uniform and colour-vision safe;
Coolwarm is the right choice for signed quantities such as charges or
potentials, where the zero point should read as neutral. Rainbow is
included because reviewers sometimes ask for it, not because it is a good
default.
\end{caltip}

Colour mapping is recomputed for every frame during trajectory playback, so
CN, GCN and property maps stay live throughout an animation.

\subsection{Bonds and vectors}

\begin{description}
  \item[\ui{Gradient bond coloring}] Blends each bond smoothly from one
    atom's colour to the other's, instead of the classic half-and-half
    split.
  \item[\ui{Force vectors}, \ui{Velocity vectors}] Draw per-atom vectors as
    lit 3D arrows. Each is enabled only when the structure actually carries
    that field; the tooltip explains what is missing when it does not.
  \item[\ui{Vector scale}] Arrow length in \AA{} per field unit,
    0.05--200, default 10.
\end{description}

\subsection{Per-element styling}

\ui{Element Settings\ldots} opens a table of the elements present in the
structure. Each row offers a colour swatch and a radius scale
(10--300\,\%, where exactly 100\,\% means "no override"), plus a
\ui{Reset} button. \ui{Reset All} clears every override at once.

\ui{Background} sets the viewport background colour. Distance fog, if
enabled, automatically fades into whatever you choose.

\section{Unit cell and axes}

The \ui{Unit Cell \& Axes} panel (zone 12):

\begin{description}
  \item[\ui{Show unit cell}] Also on \menu{View}\menusep\menu{Overlays}.
  \item[\ui{Cell color}] Wireframe colour.
  \item[\ui{Cell line width}] 1.0--8.0, default 2.0. A width of 1 draws
    plain GL lines; anything larger renders the edges as lit tubes.
  \item[\ui{Show axes triad}] The corner orientation indicator.
  \item[\ui{Axes style}] \ui{Cartesian (X, Y, Z)} or
    \ui{Lattice vectors (a1, a2, a3)}.
  \item[\ui{Axes size}] On-screen size of the triad, 48--240\,px.
\end{description}

\begin{calnote}
Core-profile OpenGL clamps hardware line width to one pixel, which is why
thick cell edges are drawn as tube geometry rather than wide lines. The
tubes are lit like the rest of the scene, so they also read correctly in
ray-traced exports.
\end{calnote}

\section{Lighting}

The \ui{Light} tab of the \ui{Visual Effects} panel (zone 9) edits up to
four directional lights with
Blinn-Phong shading. The default studio setup is three lights: a warm key
carrying the speculars, a cooler fill from the opposite side, and a back
(rim) light that separates atom silhouettes from the background.

Select a light in the list, then edit:

\begin{description}
  \item[\ui{Direction (view space)}] Three components, $-5$ to $5$. Because
    directions are in view space, the lighting stays fixed relative to the
    viewer as you orbit --- the structure turns, the studio does not.
  \item[\ui{Ambient}, \ui{Diffuse}, \ui{Specular}] Per-light colours.
\end{description}

\ui{Add} appends a fill light from the opposite side; \ui{Remove} deletes
the selected one. At least one light always remains.
